% SEGMENT OLP-0002-S01
% open-logic.tex
% compiles to produce complete text using open-logic-dev style

\documentclass[../include/open-logic-part]{subfiles}

\begin{document}

\clearpage

\begin{editorial}
இந்தக் கோப்பு Open Logic Project இல் உள்ள அனைத்துப் பாடப்பொருளையும்
ஏற்றுகிறது. இது போன்ற ஆசிரியர் குறிப்புகள் காட்டப்பட்டால்,
பாடப்பொருளை எவ்வாறு வழங்குவது என்பதைக் கருத்தில் கொள்ளாமல்
கோப்பு தொகுக்கப்பட்டுள்ளது என்பதை அவை குறிக்கின்றன.
இதை நீங்கள் வாசிக்க முடிந்தால், இந்த PDF ஐப் பயன்படுத்திக்
கற்பிப்பதோ கற்பதோ \emph{பரிந்துரைக்கப்படுவதில்லை}.

% SEGMENT OLP-0002-S02
கற்பித்தலுக்கோ தற்கற்றலுக்கோ மேலும் பொருத்தமான நூலை
உருவாக்குவதற்குப் பல வழிமுறைகளை Open Logic Project தருகிறது.
எடுத்துக்காட்டாக, இயல்பான அமைப்பில் எல்லாத் தருக்கச் செயலிகளும்
அடிப்படையானவையாகக் கொள்ளப்படுகின்றன; நிறுவல்களில் எல்லாச்
செயலிகளுக்குமான எல்லா நிகழ்வுகளும் செய்து காட்டப்படுகின்றன.
ஆனால் இவற்றில் சில நிகழ்வுகளைப் பயிற்சிகளாக விடுவது மிகவும்
சிறந்தது. Open Logic Project தொடர்ந்து உருவாக்கப்பட்டுவரும்
பணியும் ஆகும். கூட்டுப் பணியையும் மேம்பாட்டையும் ஊக்குவிக்க,
வரைவு நிலையில் உள்ள பகுதிகளும் பயிற்சிகள் இல்லாத பகுதிகளும்
மற்றும் இதுபோன்றவையும் சேர்க்கப்படுகின்றன. ஒரு பாடத்திட்டத்திற்காக
உருவாக்கப்படும் PDF இல் இத்தகைய பிரிவுகள் விலக்கப்படும்.

கற்பித்தலுக்கும் கற்றலுக்கும் மேலும் பொருத்தமான PDF நூல்களைக்
காண, \href{http://builds.openlogicproject.org/}{OLP இணையதளத்தில் உள்ள மாதிரிப் பாடத்திட்டங்களைப்}
பார்க்கவும். உங்கள் சொந்த நூலை உருவாக்க,
\href{https://github.com/OpenLogicProject/OpenLogic/tree/master/courses/sample}{மாதிரி இயக்கிக் கோப்பிலிருந்து}
தொடங்கலாம்; அல்லது மேலும் விரிவான, மேம்பட்ட எடுத்துக்காட்டுகளுக்காக,
இத்திட்டத்திலிருந்து உருவாக்கப்பட்ட பாடநூல்களின் மூலங்களைப் பார்க்கலாம்.
\end{editorial}

% SEGMENT OLP-0002-S03
\olimport[sets-functions-relations]{sets-functions-relations-complete}

\olimport[propositional-logic]{propositional-logic}

\olimport[first-order-logic]{first-order-logic}

\olimport[model-theory]{model-theory}

\olimport[computability]{computability}

\olimport[turing-machines]{turing-machines}

\olimport[incompleteness]{incompleteness}

\olimport[second-order-logic]{second-order-logic}

\olimport[lambda-calculus]{lambda-calculus}

\olimport[many-valued-logic]{many-valued-logic}

\olimport[normal-modal-logic]{normal-modal-logic}

\olimport[applied-modal-logic]{applied-modal-logic}

\olimport[intuitionistic-logic]{intuitionistic-logic}

\olimport[counterfactuals]{counterfactuals}

\olimport[set-theory]{set-theory}

\olimport[methods]{methods}

\olimport[history]{history}

\olimport[reference]{reference}

\end{document}
